\subsection{极大序域}

定义 4.——序域 K 是极大的，如果 K 的每一个序代数扩张都是平凡的。

例。--- * 我们稍后将看到（一般拓扑学，VIII，第 1 页）实数域 R 是一个极大序域。 *

极大序域的存在性是以下定理的一个推论：

定理 2.——每个序域 K 都容许一个序代数扩张，该扩张是一个极大序域。

可以证明，这个序扩张在 K-同构意义下是唯一的（VI，第 40 页，习题 15）。

设 \(\Omega\) 为 \(K\) 的一个代数闭包，并设 \(\mathfrak{N}\) 为由二元组 \((A, \omega)\) 组成的集合，其中 \(A\) 是 \(K\) 的一个扩域，包含于 \(\Omega\) ，而 \(\omega\) 是 \(A\) 上的一个序，它使 \(A\) 成为 \(K\) 的一个序扩张。给 \(\mathfrak{N}\) 定序，所用关系为 « \(L\) 是 \(M\) 的序扩张» （介于 \(M\) 与 \(L\) 之间）。配备此序后， \(\mathfrak{N}\) 是一个归纳有序集：事实上，若 \((L_{t})\) 是 \(\mathfrak{N}\) 中元素的一个全序族，则域 \(L = \cup_{t} L_{t}\) ，取 \(L_{+} = \cup_{t} (L_{t})_{+}\) 定序后，是

\(L_{i}\) 的一个上界。于是根据集合论，III，第 154 页，定理 2，N 有一个极大元，这就完成了证明。

命题 5. --- 设 K 是一个极大序域，并设 f 是 K{[}X{]} 中的多项式，它在 K 的两个元素 a 和 b（其中 a \textless{} b）之间变号。那么 f 在 K 中有一个根 x，使得 a \textless{} x \textless{} b。

f 的至少一个不可约因子，设为 h，在 a 与 b 之间变号。于是域 \(\mathrm{K}[\mathrm{X}]/(h)\) 容许 K 的一个序扩张的结构（VI，第 23 页，命题 3），且 h 的次数为 1（定义 4）。由于 \(h(a)h(b)<0\) ，h 的唯一根 x 满足 a\textless x\textless b，因为一次多项式函数是单调的。

命题 6. --- 极大序域 K 的每个正元素在 K 中都有平方根。K{[}X{]} 中每个奇数次多项式在 K 中至少有一个根。

这直接由 VI 第 24 页命题 4 的推论 2 和 1 得出。

推论。--- 在极大序域 K 上，只存在一种与域结构相容的序。

确实，\(\mathbf{K}\) 的正元素由其代数结构决定：它们就是平方数。

